\section{例题速看}
\subsection{打酱油}
\textbf{真题 · 第11次CCF-CSP认证，2017-09，第1题。}\href{https://www.acwing.com/problem/content/3250/}{打酱油（在线做题）}。

\noindent\textbf{题目与输入输出。}带N元买每瓶10元的酱油，可选择买3瓶赠1瓶或买5瓶赠2瓶，促销组之间不重复使用已赠送的瓶。求最多得到多少瓶。输入N，输出最大瓶数。N是10的整数倍且不超过300。

\noindent\textbf{读题思路。}先把钱换成可付费的瓶数，赠品不消耗钱，也不能把同一份购买重复参加两种促销。付费瓶数至多30，枚举促销组完全够用。

\noindent\textbf{做题思路与解答。}枚举买5赠2的组数f。剩余付费瓶数为r时，尽可能组成买3赠1的组，剩下的按单瓶购买，得到$7f+4\lfloor r/3\rfloor+r\bmod3$。每种方案都有一个f；固定f时多组成三瓶组只会增加赠品，所以逐项取最大值必然包含最优方案。也可证明优先五瓶组的贪心，但本题的小规模枚举不需要额外假设。
\begin{verbatim}
输入
40
输出
5
\end{verbatim}
\begin{verbatim}
输入
80
输出
11
\end{verbatim}

\noindent\textbf{复杂度与边界。}时间$O(N/50+1)$，额外空间$O(1)$。0元时答案0；余下4个付费瓶应组成三瓶组加单瓶。
\Code{49_soy.cpp}
\subsection{跳一跳}
\textbf{真题 · 第13次CCF-CSP认证，2018-03，第1题。}\href{https://www.acwing.com/problem/content/3260/}{跳一跳（在线做题）}。

\noindent\textbf{题目与输入输出。}输入若干个0、1、2，数量不超过30；0仅在末尾出现，表示游戏结束。1表示跳到方块非中心，得1分；2表示跳到中心，若是第一次跳或上次得1分，则得2分，否则比上次多2分。输出总得分。

\noindent\textbf{读题思路。}输入没有n，由0终止；连续中心得分为2、4、6等，遇到1后必须重置。末尾0不计分。

\noindent\textbf{做题思路与解答。}last表示上一跳实际得分，初值0。读到1令last=1；读到2且last不超过1时令last=2，否则加2。然后累加last。每次更新恰好实现规则，因此处理完前缀后总分始终正确。
\begin{verbatim}
输入
1 1 2 2 2 1 1 2 2 0
输出
22
\end{verbatim}

\noindent\textbf{复杂度与边界。}时间与输入数量成正比，空间$O(1)$。输入首项为0时总分为0；不要把last当作上一次输入的1或2。
\Code{50_jump.cpp}
\subsection{线性分类器}
\textbf{真题 · 第19次CCF-CSP认证，2020-06，第1题。}\csp{CSP202006A}{线性分类器}。

\noindent\textbf{题目与输入输出。}给出n个点，每点含整数坐标x、y以及类别A或B，两类都存在。随后给m条直线$\theta_0+\theta_1x+\theta_2y=0$；保证没有给定点恰在直线上。若同类点全部位于同一侧、异类点位于另一侧，输出Yes，否则No，每条直线一行。输入首行n、m，随后n行点、m行三个系数。$1\le n\le1000$，$1\le m\le20$，坐标和系数绝对值不超过$10^4$。

\noindent\textbf{评测限制。}曙梦OJ当前时限1秒，内存512MB。

\noindent\textbf{读题思路。}必须把一条线与所有点比较；不是检查两类各找一个代表就结束。Yes和No的大小写要精确。两侧由表达式符号区分，无需求斜率，因而竖直线也适用。

\noindent\textbf{做题思路与解答。}用第一个点确定“正侧应该是哪一类”，再检查每个点的符号是否与类别相符。如果有反例，这条线不满足。代入式使用long long计算；原范围下int也能容纳，但显式提升类型便于保持安全的乘法习惯。正确性来自直线两侧的表达式符号相反，以及逐点检查完整覆盖了分类条件。
\begin{verbatim}
输入
9 3
1 1 A
1 0 A
1 -1 A
2 2 B
2 3 B
0 1 A
3 1 B
1 3 B
2 0 A
0 2 -3
-3 0 2
-3 1 1
输出
No
No
Yes
\end{verbatim}

\noindent\textbf{复杂度与边界。}时间$O(nm)$，空间$O(n)$。结构体记录三项属性；读类别用格式串中的空格跳过之前的换行。
\Code{38_classifier.cpp}
\subsection{检测点查询}
\textbf{真题 · 第20次CCF-CSP认证，2020-09，第1题。}\href{https://www.acwing.com/problem/content/3295/}{检测点查询（在线做题）}。

\noindent\textbf{题目与输入输出。}输入n、居民坐标X、Y，随后n行检测点坐标，编号按输入顺序从1开始。输出最近的三个点编号，每个一行；同距离时编号小者在前。$3\le n\le200$，坐标绝对值不超过1000。

\noindent\textbf{读题思路。}要按距离、编号两层比较。无需真的开方，平方距离与距离具有相同顺序。只取三个点，三轮扫描就够用。

\noindent\textbf{做题思路与解答。}计算各点平方距离。每轮从未选点中找最小距离者，按编号升序扫描且只在距离严格减小时替换，就自动保留同距离最小编号。输出并标记已选，重复三轮。归纳可知每轮得到剩余点中的第一名，三轮恰好得到所需前三名。
\begin{verbatim}
输入
3 2 2
2 2
2 3
2 4
输出
1
2
3
\end{verbatim}
\begin{verbatim}
输入
5 0 1
-1 0
0 0
1 0
0 2
-1 2
输出
2
4
1
\end{verbatim}

\noindent\textbf{复杂度与边界。}时间$O(3n)=O(n)$，数组空间$O(n)$。平方距离最大为$8\times10^6$，此范围int可装下；代码用long long并先提升差值类型。
\Code{51_nearest.cpp}
\subsection{田地丈量}
\textbf{真题 · 第29次CCF-CSP认证，2023-03，第1题。}\csp{CSP202303A}{田地丈量}。

\noindent\textbf{题目与输入输出。}输入n、a、b，随后n行现有田地矩形的左下与右上坐标。现有田地互不重叠（交集面积为0），目标区域为$[0,a]\times[0,b]$。求目标区域内已有田地面积。$n\le100$，所有坐标绝对值不超过$10^4$，a、b为正。输出面积。

\noindent\textbf{读题思路。}统计连续面积，边接触贡献0，不能加1。现有田地互不重叠，所以各块与目标的交面积可以直接相加；若取消这一保证，则不能简单累加。

\noindent\textbf{做题思路与解答。}每块田地裁剪到目标：横向左端取$\max(x_1,0)$、右端取$\min(x_2,a)$，纵向同理。有正宽、正高才累加其乘积。两个一维交区间的笛卡尔积正好是矩形交集；原田地互不重叠，因此求和无重复覆盖。
\begin{verbatim}
输入
4 10 10
0 0 5 5
5 -2 15 3
8 8 15 15
-2 10 3 15
输出
44
\end{verbatim}

\noindent\textbf{复杂度与边界。}时间$O(n)$，额外空间$O(1)$，可边读边算。面积与累加值用long long，乘法前提升类型。
\Code{52_land.cpp}
\subsection{小中大}
\textbf{真题 · 第16次CCF-CSP认证，2019-03，第1题。}\href{https://www.acwing.com/problem/content/3275/}{小中大（在线做题）}。

\noindent\textbf{题目与输入输出。}输入$n\le10^5$以及n个已按升序或降序排列的整数，可重复，绝对值不超过$10^7$。输出最大值、中位数、最小值，空格分隔；整数不带小数，非整数保留1位小数。

\noindent\textbf{读题思路。}已经有序，无需重排；升序和降序都能由首尾取得极值。奇数长度中位数是一项，偶数长度是中间两项均值。

\noindent\textbf{做题思路与解答。}首尾取min、max，中间位置按奇偶取值。先计算中位数的两倍twice；偶数时直接输出twice/2，否则输出twice/2.0保留一位。这使负半整数也正确，避免整数除法丢掉0.5。
\begin{verbatim}
输入
3
-1 2 4
输出
4 2 -1
\end{verbatim}
\begin{verbatim}
输入
4
-2 -1 3 4
输出
4 1 -2
\end{verbatim}

\noindent\textbf{复杂度与边界。}读入时间$O(n)$、数组空间$O(n)$。可以只保留首尾及中间项把额外空间降为$O(1)$。n=1时三个输出值相同。
\Code{53_small_middle_large.cpp}
\subsection{数组推导}
\textbf{真题 · 第23次CCF-CSP认证，2021-09，第1题。}\href{https://www.acwing.com/problem/content/4009/}{数组推导（在线做题）}。

\noindent\textbf{题目与输入输出。}非负整数数组A的前缀最大值为$B_i=\max(A_1,\ldots,A_i)$。已知单调不降的B，求A的元素和最大值及最小值，各输出一行，最大值在前。输入n、n项B；$1\le n\le100$，$0\le B_i\le10^5$。

\noindent\textbf{读题思路。}B是前缀最大值，不是前缀和。B上升时当前A被迫等于新值；B不变时A可小到0。

\noindent\textbf{做题思路与解答。}最大和为$\sum B_i$，因为各项都可选为B对应值，且不能超过该值。最小和只保留第一项及严格上升处的B值，其他A取0。上升时此前元素不可能提供新最大值，必须由当前A提供；这些强制项已足以实现全部B，所以该和最小。
\begin{verbatim}
输入
6
0 0 5 5 10 10
输出
30
15
\end{verbatim}
\begin{verbatim}
输入
7
10 20 30 40 50 60 75
输出
285
285
\end{verbatim}

\noindent\textbf{复杂度与边界。}时间$O(n)$，额外空间$O(1)$。B首项可为0，全0时两个答案都是0。
\Code{54_derive.cpp}
\subsection{如此编码}
\textbf{真题 · 第27次CCF-CSP认证，2022-09，第1题。}\href{https://www.acwing.com/problem/content/4702/}{如此编码（在线做题）}。

\noindent\textbf{题目与输入输出。}每题有$a_i\ge2$个选项，正确选项$0\le b_i<a_i$。设$c_0=1,c_i=\prod_{j=1}^i a_j$，编码$m=\sum_{i=1}^n c_{i-1}b_i$。输入n、m及n项a，恢复并输出n项b，空格分隔。$n\le20,c_n\le10^9,0\le m<c_n$。

\noindent\textbf{读题思路。}这是不同位使用不同进制的编码。题面给出长公式，但可以从最低位逐次取余，避免反复计算前缀积。

\noindent\textbf{做题思路与解答。}对$a_1$取余，后面所有项都是$a_1$的倍数，所以$b_1=m\bmod a_1$。再令$m=\lfloor m/a_1\rfloor$，得到余下位的同类编码。依次取余和整除即可恢复全部b。每步都唯一确定一位，最终得到原编码的唯一表示。
\begin{verbatim}
输入
4 0
2 3 2 5
输出
0 0 0 0
\end{verbatim}
\begin{verbatim}
输入
7 23333
3 5 20 10 4 3 10
输出
2 2 15 7 3 1 0
\end{verbatim}

\noindent\textbf{复杂度与边界。}时间$O(n)$，额外空间$O(1)$，可边读a边输出。m为0时每位都为0；余数不是输出的逆序。
\Code{55_encoding.cpp}
\subsection{相反数与相邻数对：不重不漏}

\noindent\textbf{做题：}\csp{CSP201403A}{相反数（真题，第1次，2014-03，第1题）}、\csp{CSP201409A}{相邻数对（真题，第2次，2014-09，第1题）}。

\noindent\textbf{解答。}这两题数据小，读入数组后枚举$i<j$即可。相反数检查$a_i+a_j=0$；相邻数对检查差为1或-1。每个无序对只枚举一次。题面保证输入值互不相同；若改成可重复，按值存在性计数便不等于按位置数对计数。
\begin{lstlisting}
int answer = 0;
for (int i = 0; i < n; i++) {
    for (int j = i+1; j < n; j++) {
        if (1LL*a[i]+a[j] == 0) answer++; // Opposite numbers.
        // Neighbor pairs: replace the condition above by
        // 1LL*a[i]-a[j] == 1 || 1LL*a[i]-a[j] == -1
    }
}
\end{lstlisting}
时间$O(n^2)$，数组空间$O(n)$。值域小且不重复时也可用存在性桶：相反数只检查正数x与-x，相邻数对只检查x与x+1，这同样只计一次。
\subsection{门禁系统与灰度直方图：同一种桶的两种输出}

\noindent\textbf{做题：}\csp{CSP201412A}{门禁系统（真题，第3次，2014-12，第1题）}、\csp{CSP202104A}{灰度直方图（真题，第22次，2021-04，第1题）}。
门禁题输入n及n次记录，每条是编号，输出每条编号在当前前缀中的出现次数；原题$1\le n\le1000$，编号在1到n之间。直方图题输入n、m、L和$n\times m$个灰度值，灰度在0到$L-1$之间，输出L个计数。

\noindent\textbf{解答。}数组cnt全0，每读入一个x就\texttt{cnt[x]++}。门禁在更新后立即输出，直方图在全部更新后按0到L-1输出。
\begin{lstlisting}
int cnt[1005] = {0};
for (int i = 0; i < n; i++) {
    int id; scanf("%d", &id);
    printf("%d%c", ++cnt[id], i+1 == n ? '\n' : ' ');
}
// Histogram core: read all pixels and do cnt[gray]++.
// Then print cnt[0], cnt[1], ..., cnt[L-1], including zeros.
\end{lstlisting}
每条记录只访问一次，时间线性；存储取决于值域，不是记录数量。多个测试实例时必须重新清零。
\subsection{数列分段与分蛋糕：何时增加答案}

\noindent\textbf{做题：}\csp{CSP201509A}{数列分段（真题，第5次，2015-09，第1题）}、\csp{CSP201703A}{分蛋糕（真题，第10次，2017-03，第1题）}。
数列分段把相邻相等项放在同一段，输入n及数组，输出段数；n至少为1。

\noindent\textbf{解答。}先有一段，每个$a_i\ne a_{i-1}$的位置开启新段，答案为1加变化次数。
\begin{lstlisting}
int segments = 1;
for (int i = 1; i < n; i++) if (a[i] != a[i-1]) segments++;
\end{lstlisting}
分蛋糕按输入次序累积正重量，达到k就交给一个人并从0重新累积；最后不足k但仍有蛋糕也交给一个人。输入n、k和各块重量，输出人数。
\begin{lstlisting}
long long sum = 0; int people = 0;
for (int i = 0; i < n; i++) {
    sum += weight[i];
    if (sum >= k) { people++; sum = 0; }
}
if (sum > 0) people++;
\end{lstlisting}
达到阈值就分出，符合题目指定流程；若从最大化可分人数理解，正重量条件下延后分出只会消耗更多块，不能增加后续可分段数。两题时间$O(n)$；读入可流式处理，额外空间$O(1)$。
\subsection{卖菜：首尾与同步更新}

\noindent\textbf{做题：}\csp{CSP201809A}{卖菜（真题，第14次，2018-09，第1题）}。输入至少2个原价格；内部新价格是左、中、右原价格之和除以3向下取整，两端是两个相邻原价格之和除以2。输出新价格序列。

\noindent\textbf{解答与关键代码。}全部原值先读入a，再按规则输出，或者写入新数组b；不要修改a后再算后面的均值。价格非负，整数除法正好完成向下取整。
\begin{lstlisting}
printf("%d", (a[0]+a[1])/2);
for (int i = 1; i+1 < n; i++) printf(" %d", (a[i-1]+a[i]+a[i+1])/3);
printf(" %d\n", (a[n-2]+a[n-1])/2);
\end{lstlisting}
时间$O(n)$，空间$O(n)$。n=2时内部循环不执行，两个答案都是两项均值。
\subsection{报数：数字与轮次分开}

\noindent\textbf{做题：}\csp{CSP201912A}{报数（真题，第18次，2019-12，第1题）}。四人依次轮流报自然数，从1开始；7的倍数或十进制含7时跳过并记录当前人的跳过次数。直到成功报出n个数，输出四人的跳过次数（每人一行）。

\noindent\textbf{解答。}变量x记录当前自然数，done只统计成功次数；每处理一个x，轮次都会推进，包括被跳过的数。当前人是$(x-1)\bmod4$。
\begin{lstlisting}
bool skip(int x) {
    if (x % 7 == 0) return true;
    while (x) { if (x % 10 == 7) return true; x /= 10; }
    return false;
}
// int cnt[4]={0}, done=0;
// for (int x=1; done<n; x++)
//     if (skip(x)) cnt[(x-1)%4]++; else done++;
\end{lstlisting}
不要在跳过时停留在同一个人，也不要把x直接限制为n。若实际处理T个数，时间$O(T\log T)$，空间$O(1)$。
\subsection{归一化处理：均值、方差与浮点除法}

\noindent\textbf{做题：}\csp{CSP202206A}{归一化处理（真题，第26次，2022-06，第1题）}。给定n个整数，计算均值$\mu=\frac1n\sum a_i$、总体方差$D=\frac1n\sum(a_i-\mu)^2$，再逐行输出$(a_i-\mu)/\sqrt D$；原题保证方差非零。

\noindent\textbf{解答与关键代码。}使用double累加，先求均值，再求平方偏差和。不要误用样本方差的分母n-1，也不要把平方根漏掉。
\begin{lstlisting}
double sum = 0, variance = 0;
for (int i = 0; i < n; i++) sum += a[i];
double mean = sum/n;
for (int i = 0; i < n; i++) {
    double d = a[i]-mean; variance += d*d;
}
variance /= n;
for (int i = 0; i < n; i++) printf("%.10f\n", (a[i]-mean)/sqrt(variance));
\end{lstlisting}
需包含\texttt{<cmath>}并调用相应sqrt。时间$O(n)$，数组空间$O(n)$。自编变式若允许全部相等，必须规定D=0时的输出，不能直接除以0。
\subsection{重复局面：不用STL也能比较字符串}

\noindent\textbf{做题：}\csp{CSP202305A}{重复局面（真题，第30次，2023-05，第1题）}。输入n个$8\times8$字符棋盘，按输入顺序输出当前棋盘在截至此刻的全部棋盘中出现了几次。原题$n\le100$。

\noindent\textbf{解答。}每个棋盘压成一个长度64的字符数组，末尾加0；枚举此前及当前所有棋盘，用strcmp判断完全相同。n很小，$O(64n^2)$足够，
\begin{lstlisting}
char board[105][65];
// Read each row into char row[9] with scanf("%8s", row),
// then copy its 8 characters to board[i]+8*r.
// Set board[i][64]=0 before strcmp.
for (int i = 0; i < n; i++) {
    int count = 0;
    for (int j = 0; j <= i; j++)
        if (strcmp(board[i], board[j]) == 0) count++;
    printf("%d\n", count);
}
\end{lstlisting}
包含\texttt{<cstring>}。比较必须覆盖全部64格，不能只用某一个格子或棋子数量作为身份。

\subsection{图像旋转}
\textbf{真题 · 第4次CCF-CSP认证，2015-03，第1题。}\csp{CSP201503A}{图像旋转}。

\noindent\textbf{题目与输入输出。}输入一个$n\times m$矩阵，将它逆时针旋转$90^\circ$后输出。第一行n、m，随后n行每行m个整数；输出m行，每行n个整数。$1\le n,m\le1000$，元素在0到1000之间。

\noindent\textbf{评测限制。}曙梦OJ当前时限1秒，内存256MB。

\noindent\textbf{读题思路。}圈出“逆时针”和“输出m行n列”。矩阵不一定是正方形，不能把行列长度混用。百万个int约3.8MiB，可以存整个原矩阵。

\noindent\textbf{做题思路与解答。}原矩阵第j列从上到下，会成为新矩阵第$m-1-j$行（0下标）。因此按列从右到左、每列从上到下输出即可；无需再开一张新矩阵。映射式为$b[m-1-j][i]=a[i][j]$。每个输入元素恰好输出一次，所以既没有漏项，也没有重复。
\begin{verbatim}
输入
2 3
1 5 3
3 2 4
输出
3 4
5 2
1 3
\end{verbatim}

\noindent\textbf{复杂度与边界。}时间$O(nm)$，空间$O(nm)$。用$1\times3$和$3\times1$检查方向；每行最后一个值后输出换行。
\Code{35_rotate.cpp}
\subsection{数位之和}
\textbf{真题 · 第6次CCF-CSP认证，2015-12，第1题。}\csp{CSP201512A}{数位之和}。

\noindent\textbf{题目与输入输出。}输入一个整数$n$，输出其十进制各位数字之和。$0\le n\le10^9$。

\noindent\textbf{评测限制。}曙梦OJ当前时限1秒，内存256MB。

\noindent\textbf{读题思路。}任务是求数字之和，不是数位数量，也不是把各位拼成新数。输入允许0；不必为了拆位把整数转换为字符串。

\noindent\textbf{做题思路与解答。}每轮用$n\bmod10$取个位，再用整数除法$n/10$删除个位；不断累加。每位恰好取一次，得到所求总和。do-while会对0执行一次，结果仍为0；普通while配合总和初值0也正确。
\begin{verbatim}
输入
20151220
输出
13
\end{verbatim}

\noindent\textbf{复杂度与边界。}若十进制位数为d，时间$O(d)$，空间$O(1)$。$10^9$只有10位，和最多81（$10^9$本身为1）；int足够。
\Code{36_digits.cpp}
\subsection{折点计数}
\textbf{真题 · 第7次CCF-CSP认证，2016-04，第1题。}\csp{CSP201604A}{折点计数}。

\noindent\textbf{题目与输入输出。}有n天销售量$a_1,\ldots,a_n$，相邻两天的销售量不同。一个内部位置的值严格大于左右两项，或严格小于左右两项，称为折点。输入n及n项销售量，输出折点数。$1\le n\le1000$，销售量在0到10000之间。

\noindent\textbf{评测限制。}曙梦OJ当前时限1秒，内存256MB。

\noindent\textbf{读题思路。}折点表示趋势从升变降或从降变升。第一天和最后一天没有两侧邻居，不能计入；不相邻的位置可以取相同值。

\noindent\textbf{做题思路与解答。}从第二项到倒数第二项检查“同时大于左右”或“同时小于左右”。每个内部位置独立判断一次，恰好覆盖全部折点。也可检查相邻差的符号是否不同，但直接写大小关系更直观，且没有乘法溢出的隐患。
\begin{verbatim}
输入
7
5 4 1 2 3 6 4
输出
2
\end{verbatim}

\noindent\textbf{复杂度与边界。}时间$O(n)$，空间$O(n)$。n为1或2时循环不执行，答案为0。严格单调序列也没有折点。
\Code{37_turns.cpp}
\subsection{仓库规划}
\textbf{真题 · 第32次CCF-CSP认证，2023-12，第1题。}\csp{CSP202312A}{仓库规划}。

\noindent\textbf{题目与输入输出。}n个仓库各有m维坐标。只有j的每一维都严格大于i，j才是i的上级；若有多个，取编号最小的j，否则为0。输入n、m，再输入n行坐标；按编号顺序输出各仓库上级，一行一个。$1\le n\le1000$，$1\le m\le10$，各坐标绝对值不超过$10^6$。

\noindent\textbf{评测限制。}曙梦OJ当前时限1秒，内存512MB。

\noindent\textbf{读题思路。}关注“每一维”“严格大于”“编号最小”。输出不是最近仓库，也不是坐标最小仓库。$n^2m\le10^7$，直接枚举候选可行，不必急着排序或使用STL。

\noindent\textbf{做题思路与解答。}对每个i，按编号递增枚举j；逐维检查，一旦出现$a[j][k]\le a[i][k]$便拒绝。第一个全部通过的j就是编号最小答案，可以退出候选循环。自身必然无法通过严格比较。遍历所有候选保证无遗漏，升序与首次接受保证最小编号。
\begin{verbatim}
输入
4 2
0 0
-1 -1
1 2
0 -1
输出
3
1
0
3
\end{verbatim}

\noindent\textbf{复杂度与边界。}时间$O(n^2m)$，空间$O(nm)$。相等坐标不是上级；无上级时保留初值0。
\Code{39_warehouses.cpp}
\subsection{统计一组整数（教学例题）}

\noindent\textbf{题目。}输入 $n$ 个整数，输出总和、最小值、最大值，以及出现次数最多的值及其次数。若多个值并列，输出最小的值。$1\le n\le2\times10^5$，$|a_i|\le10^9$。

\noindent\textbf{输入。}第一行 $n$，第二行 $n$ 个整数。

\noindent\textbf{输出。}第一行依次输出总和、最小值、最大值；第二行输出众数及出现次数。
\begin{verbatim}
输入
6
3 -1 3 2 -1 5
输出
11 -1 5
-1 2
\end{verbatim}

\noindent\textbf{解答。}总和与极值可边读边更新；最小值不能初始化为0，否则全正数时会出错。众数作为本例的进阶要求：先调用C标准库qsort排序，再扫描每个相等值段。按值升序扫描且只在次数严格增加时更新答案，就保留了并列时较小的值。qsort的比较函数返回负数、0或正数，表示小于、等于或大于；不要用可能溢出的相减判断。扫描为$O(n)$，空间为$O(n)$，总时间由排序主导。C标准未规定qsort的复杂度保证；需要明确的$O(n\log n)$比较次数保证时，可使用第二篇讲解的sort。初学时先掌握总和与极值，再学习众数部分。
\Code{01_stats.cpp}\subsection{网格的一轮同步更新（教学例题）}

\noindent\textbf{题目。}给定一个 $n\times m$ 的0/1网格。一次更新后，每个格子的新值由其旧状态的上下左右四邻格决定：值为1的邻格至少有两个，则新值为1，否则为0。越界邻格不计入。只进行一轮更新，$1\le n,m\le500$。

\noindent\textbf{输入。}第一行 $n,m$，随后每行 $m$ 个0或1。

\noindent\textbf{输出。}更新后的网格。
\begin{verbatim}
输入
2 3
1 0 1
0 1 0
输出
0 1 0
1 0 1
\end{verbatim}

\noindent\textbf{推演。}左上格的两个有效邻格都是0，所以变为0；上方中间格的三个邻格都是1，所以变为1。所有计数都使用原数组 a，把答案放进 b。若直接覆盖 a，后面的格子会读到已经变化的数据。

\noindent\textbf{正确性。}程序对每格恰好检查四个方向，只有合法坐标参与计数，因此统计值与题意一致；写入 b 不改变 a，满足同步更新。时间和空间均为 $O(nm)$。
\Code{02_grid.cpp}
